\documentclass[a4paper]{article}
% Espanol (Mexico) con XeLaTeX: fontspec para fuentes Unicode, polyglossia
% para la division silabica y los nombres del documento en la variante mexicana.
\usepackage{fontspec}
\usepackage{polyglossia}
\setdefaultlanguage[variant=mexican]{spanish}
% polyglossia no traduce los nombres de listings.
\AtBeginDocument{\providecommand{\lstlistingname}{}\renewcommand{\lstlistingname}{Listado}%
  \providecommand{\lstlistlistingname}{}\renewcommand{\lstlistlistingname}{Índice de listados}}
\usepackage{amsmath,amssymb}
\usepackage{graphicx}
\usepackage[margin=2.35cm]{geometry}
\usepackage[most]{tcolorbox}
\usepackage{etoolbox}
\usepackage{listings}
\usepackage{xcolor}
\usepackage{hyperref}
\usepackage{booktabs,longtable,tabularx,array}
\usepackage{subcaption}
\usepackage{float}
\usepackage{enumitem}
\usepackage{microtype}
\usepackage{fancyhdr}
\usepackage{needspace}

\definecolor{kbBack}{HTML}{F2F6FA}
\definecolor{kbFrame}{HTML}{9DB4CC}
\definecolor{kbTitle}{HTML}{52708F}
\definecolor{ibBack}{HTML}{FBF5E7}
\definecolor{ibFrame}{HTML}{C9A961}
\definecolor{ibTitle}{HTML}{7D5F1E}
\definecolor{wbBack}{HTML}{FCF2F0}
\definecolor{wbFrame}{HTML}{D6A096}
\definecolor{wbTitle}{HTML}{9E4F43}
\definecolor{dbBack}{HTML}{F4F7F3}
\definecolor{dbFrame}{HTML}{AEC2AC}
\definecolor{dbTitle}{HTML}{5F7F64}

\tcbset{softbox/.style={
  enhanced,breakable,boxrule=0.8pt,arc=2pt,left=3mm,right=3mm,
  top=2mm,bottom=2mm,coltitle=white,fonttitle=\bfseries,
  toptitle=0.6mm,bottomtitle=0.6mm
}}
\newtcolorbox{knowledgebox}[1]{softbox,colback=kbBack,colframe=kbFrame,colbacktitle=kbTitle,title={#1}}
\newtcolorbox{importantbox}[1]{softbox,colback=ibBack,colframe=ibFrame,colbacktitle=ibTitle,title={#1}}
\newtcolorbox{warningbox}[1]{softbox,colback=wbBack,colframe=wbFrame,colbacktitle=wbTitle,title={#1}}
\newtcolorbox{dialoguebox}[1]{softbox,colback=dbBack,colframe=dbFrame,colbacktitle=dbTitle,title={#1}}

\lstset{
  language=Python,basicstyle=\ttfamily\small,
  keywordstyle=\color{kbTitle}\bfseries,stringstyle=\color{wbTitle},
  commentstyle=\color{dbTitle},breaklines=true,frame=single,numbers=left,
  numberstyle=\tiny\color{gray},captionpos=b,columns=fullflexible,
  keepspaces=true,showstringspaces=false,extendedchars=true
}
\BeforeBeginEnvironment{lstlisting}{\Needspace{12\baselineskip}}

\hypersetup{colorlinks=true,linkcolor=kbTitle,urlcolor=kbTitle,citecolor=wbTitle}
\setlist{itemsep=0.25em,topsep=0.4em}
\setlength{\parindent}{2em}
\setlength{\parskip}{0.25em}
\newcolumntype{L}[1]{>{\raggedright\arraybackslash}p{#1}}

\providecommand{\notetitle}{Notas del curso: [tema del curso]}
\providecommand{\noteauthors}{Elaboradas a partir de los materiales públicos de Stanford CS153}
\providecommand{\notedate}{Primavera de 2026}
\providecommand{\videochannel}{Stanford Online}
\providecommand{\videopublishdate}{2026}
\providecommand{\videoduration}{[duración del video]}
\providecommand{\videourl}{}
\providecommand{\videocoverpath}{cover.jpg}
\providecommand{\coursesite}{https://cs153.stanford.edu/}
\providecommand{\repourl}{https://github.com/hqhq1025/ai-course-notes}
\providecommand{\skillversion}{wdkns/wdkns-skills@39f1a04}

\newcommand{\figsource}[1]{\par\vspace{-0.3em}{\footnotesize\color{gray}\emph{Fuente: #1}}\par}
\newcommand{\teachervoice}[1]{\begin{knowledgebox}{Nota de clase}#1\end{knowledgebox}}
\newcommand{\term}[1]{\textbf{#1}}

\pagestyle{fancy}
\fancyhf{}
\fancyfoot[C]{\thepage}
\fancyfoot[R]{\footnotesize\color{gray}\href{\repourl}{AI Course Notes}}
\renewcommand{\headrulewidth}{0pt}
\renewcommand{\footrulewidth}{0pt}

\newcommand{\makecscover}{
\begin{titlepage}
\centering
\vspace*{0.7cm}
{\Large Stanford CS153 · Frontier Systems\par}
\vspace{0.8cm}
{\huge\bfseries \notetitle\par}
\vspace{0.6cm}
{\large \noteauthors\par}
\vspace{0.25cm}
{\large \notedate\par}
\vspace{0.9cm}
\ifdefempty{\videocoverpath}{}{
  \includegraphics[width=0.86\textwidth,height=0.43\textheight,keepaspectratio]{\videocoverpath}\par
}
\vfill
\begin{tcolorbox}[width=0.92\textwidth,colback=black!2!white,colframe=black!45,boxrule=0.7pt,arc=2pt]
\textbf{Autor o canal del video}: \videochannel\par
\textbf{Fecha de publicación}: \videopublishdate\par
\textbf{Duración del video}: \videoduration\par
\textbf{Sitio del curso}: \href{\coursesite}{\nolinkurl{\coursesite}}\par
\ifdefempty{\videourl}{}{\textbf{Enlace del video}: \href{\videourl}{\nolinkurl{\videourl}}\par}
\textbf{Normas de elaboración}: \skillversion\ + las normas source-first / teacher-voice / visual-QA de este repositorio
\end{tcolorbox}
\end{titlepage}
\tableofcontents
\newpage
}


\renewcommand{\notetitle}{Lecture 09: Identity Infrastructure---Trust, Incident Learning y Agent Delegation}
\renewcommand{\noteauthors}{Elaborado a partir de la entrevista a Todd McKinnon, los subtítulos con marcas de tiempo y materiales de Okta/Auth0 y de estándares abiertos}
\renewcommand{\notedate}{Curso de invierno de 2025 · versión reescrita en 2026 con enfoque source-first}
\renewcommand{\videochannel}{Grabaciones de ediciones anteriores de Stanford CS153}
\renewcommand{\videopublishdate}{2025-02-11}
\renewcommand{\videoduration}{39:02}
\renewcommand{\videourl}{https://www.youtube.com/watch?v=wu2BWTVQQ1Q}
\renewcommand{\videocoverpath}{cover.jpg}

\newcommand{\lecturefigure}[3]{%
\begin{figure}[H]
  \centering
  \includegraphics[width=0.96\textwidth]{images/#1}
  \caption{#2}
  \figsource{#3}
\end{figure}
}

\begin{document}
\makecscover

\section{Auditoría de fuentes: convertir una entrevista con un CEO en un curso de Identity Systems}

Esta clase la imparte Todd McKinnon, cofundador y CEO de Okta. La sesión recorre la ventana para emprender en la computación en la nube de 2009, la infraestructura de identidad, los incidentes de seguridad, la cultura de la empresa, la adquisición de Auth0 y los AI agent. El repositorio conserva la portada oficial y 1,033 subtítulos con marca de tiempo. Este texto toma como eje de la clase la entrevista de febrero de 2025 y verifica los durable mechanism con materiales oficiales de Okta/Auth0 y con los estándares de NIST, IETF, OpenID y W3C. No presenta las capacidades de producto de 2026, como non-human identity, continuous session risk o Secure AI, como si ya estuvieran disponibles en el momento de la clase.

\begin{warningbox}{Límites de la evidencia sobre las cifras, los incidentes y los juicios organizacionales de la clase}
Las cifras de usuarios que inician sesión al mes, el volumen de solicitudes de autenticación, el número de empleados, la situación de los clientes y los juicios sobre el mercado que menciona el ponente se conservan \textbf{tal como se dijeron en clase} y solo sirven para ilustrar la escala de la workload y de la responsabilidad. La parte de incidentes de seguridad verifica los mecanismos de credential, session, logging y remediation con los RCA públicos, pero un incidente complejo no puede atribuirse a un solo empleado. Las opiniones sobre la adquisición y sobre la IA son juicios basados en la experiencia de McKinnon, no leyes causales universales.
\end{warningbox}

\begin{importantbox}{Los tres ciclos cerrados de esta clase}
\begin{enumerate}
  \item \textbf{Access loop}: identity source $\rightarrow$ policy $\rightarrow$ authentication/authorization $\rightarrow$ session/token $\rightarrow$ audit;
  \item \textbf{Trust loop}: incident $\rightarrow$ containment $\rightarrow$ root cause $\rightarrow$ remediation/notification $\rightarrow$ customer evidence;
  \item \textbf{Transition loop}: technology shift $\rightarrow$ focused wedge $\rightarrow$ infrastructure depth $\rightarrow$ organizational adaptation.
\end{enumerate}
\end{importantbox}

\begin{knowledgebox}{Términos clave: Identity no es una «tabla de cuentas»}
\term{principal} (sujeto) es una entidad que puede autenticarse y autorizarse; puede ser una persona, una service account, una workload o un agent. El \term{identity provider} (proveedor de identidad, IdP) mantiene el estado de los sujetos y proporciona a las aplicaciones aserciones de identidad confiables. La \term{authentication} (autenticación) responde «¿quién eres?»; la \term{authorization} (autorización) responde «¿qué puedes hacer con qué recurso?». La \term{session/token}, por su parte, convierte el resultado de una verificación en una capacidad de acceso reutilizable durante un tiempo limitado.
\end{knowledgebox}

\subsection{Por qué la página de inicio de sesión es solo la punta del iceberg}

Esta sección define primero el objeto de estudio del curso. Lo que el usuario ve es un login redirect o un MFA challenge; lo que la plataforma mantiene en realidad es directory, policy, protocol, session, provisioning, availability, security monitoring y audit. El servicio de identidad está en la «puerta de entrada» de muchas aplicaciones: un allow erróneo puede convertirse en una filtración de datos, y un deny erróneo o un outage puede impedir que toda la organización trabaje. Por eso pertenece a la vez a la infraestructura security-critical y availability-critical.

\subsection{Resumen de la sección}

El protagonista de la Lecture 09 no es la historia de emprendimiento de una empresa, sino el identity control plane. Todas las experiencias organizacionales que siguen se traducirán en preguntas de sistemas: dónde está el estado, cómo se toman las decisiones, cómo se acotan las fallas y cómo se recupera la confianza con evidencia.
\section{Transición tecnológica y Wedge: por qué Cloud Identity}

McKinnon entiende la oportunidad de emprender dentro de la cloud transition: cuando las aplicaciones empresariales pasan de centros de datos propios a SaaS, el antiguo perímetro de red y los directorios locales ya no cubren de forma natural todos los recursos. Una empresa nueva no puede replicar de frente el mercado antiguo en el que el incumbent es más fuerte; debe encontrar el nuevo control point que la transición deja expuesto y entrar con un wedge por el que el cliente esté dispuesto a pagar de inmediato.

\subsection{Del Monitoring Pivot a SSO}

El equipo empezó con cloud application monitoring, pero el buyer feedback fue tibio; en cambio, ante varias cuentas de SaaS como Gmail y Dropbox, los inicios de sesión repetidos, las contraseñas y los permisos de quienes dejan la empresa se volvieron un problema inmediato para los empleados. \term{single sign-on} (inicio de sesión único, SSO) permite al usuario acceder a varias aplicaciones a través de un IdP de confianza, lo que reduce las credenciales repetidas y centraliza la policy. SSO es solo un wedge, pero por su naturaleza exige conectarse con muchas aplicaciones, directorios y protocolos.

\lecturefigure{01-transition-window.png}{La transición tecnológica crea nuevos puntos de control; el identity wedge crece desde el problema inmediato del inicio de sesión hasta convertirse en infraestructura compartida.}{Subtítulos locales 00:00--05:20; redibujo conceptual.}

\begin{knowledgebox}{Lectura de la figura: Transition, pain y platform deben cumplirse a la vez}
La cloud transition cambia dónde están las aplicaciones y cómo se compran; la credential fragmentation genera un pain concreto; SSO ofrece un resultado de corto plazo por el que se paga; directory, policy y lifecycle dan al producto profundidad de plataforma. Con transition pero sin buyer pain, se convierte en una investigación de largo plazo; con un problema pequeño pero sin superficie de expansión, se queda en feature; con platform vision pero sin wedge, es difícil obtener el primer tráfico de producción.
\end{knowledgebox}

\teachervoice{McKinnon no idealiza la primera idea: una vez construido el monitoring prototype, los clientes potenciales no mostraron entusiasmo, y por eso el equipo giró hacia «hacer más simple el inicio de sesión en SaaS». La nota de clase: primero buscar una willingness to pay real y después demostrar que detrás de la pequeña entrada existe una infrastructure con suficiente profundidad.}

\subsection{El iceberg de la infraestructura: la complejidad detrás de una experiencia simple}

Que un empleado haga clic en el ícono de una aplicación parece solo «escribir la contraseña una vez menos»; sin embargo, el sistema debe sincronizar el estado del usuario, elegir la fuente de identidad, aplicar las políticas de autenticación, emitir la assertion/token, establecer la session, gestionar logout y revoke, y actualizar los entitlements cuando la organización cambia. El efecto acumulativo del identity wedge proviene de que estos primitive compartidos pueden atender cada vez más aplicaciones, no de la página de inicio de sesión en sí.

\begin{importantbox}{Las dos restricciones del Wedge}
Una entrada de emprendimiento en infraestructura necesita a la vez: \textbf{today value}, que el cliente esté dispuesto a pagar hoy por menos fricción al iniciar sesión, un onboarding más rápido o un offboarding más controlable; y \textbf{tomorrow depth}, que la entrada pueda extenderse a directory, MFA, policy, provisioning, governance y threat response. Si falta cualquiera de los dos, el producto se queda sin flujo de efectivo o sin espacio de plataforma a largo plazo.
\end{importantbox}

\subsection{Resumen de la sección}

La oportunidad de cloud identity proviene del cambio en los límites de las aplicaciones. SSO es la puerta de entrada que se puede vender; identity lifecycle y policy son la plataforma posterior. El Pivot se fundamenta en la evidencia de los clientes, no en el costo hundido del fundador en su primer prototype.
\section{Identity Control Plane: State, Policy y Lifecycle}

La sección anterior explicó el wedge; esta entra en el sistema propiamente dicho. El identity control plane reúne HR, directory, customer signup, partner federation y service identity, que están dispersos, en un principal model unificado, y después aplica a cada acceso las políticas de authentication, authorization, session y lifecycle. No es la única truth de todos los datos, pero debe saber qué source tiene autoridad sobre qué attribute.

\subsection{El directory es la vista unificada del estado de identidad}

El \term{directory} (directorio) guarda el principal profile, los group, las relationship, el credential enrollment y el lifecycle state. La durable idea de Okta Universal Directory consiste en mapear las identity heterogéneas de Active Directory, LDAP, HR, las aplicaciones y las API a un modelo de datos predecible. La \term{source of truth} (fuente de autoridad) debe definirse por campo: por ejemplo, HR decide el employment status, la aplicación decide el role local y el IdP decide el authenticator enrollment; no se puede usar «el último en escribir» para sobrescribir toda la semántica.

\lecturefigure{02-identity-control-plane.png}{El identity control plane conecta las fuentes de identidad, el directorio unificado, el motor de políticas, las aplicaciones y la auditoría.}{Subtítulos locales 03:00--06:10; concepto oficial de Okta Universal Directory; redibujo conceptual.}

\begin{knowledgebox}{Lectura de la figura: primero pregunta de dónde viene el estado y adónde va}
La identity source se encarga de crear o actualizar los hechos; el directory normaliza principal, group y attribute; el policy engine toma decisiones de acceso según la aplicación de destino y el contexto; las apps y el audit reciben token, provisioning event y registros del sistema. Lo verdaderamente difícil son el conflict, la freshness, la deletion y la ownership: quién prevalece cuando HR y el directorio local están en conflicto, cuánto tarda en propagarse un evento de baja, cómo se recalcula un derived group y cómo demuestra la auditoría qué versión del estado usó un allow determinado.
\end{knowledgebox}

\subsection{Authentication, authorization, provisioning y audit}

La \term{multi-factor authentication} (autenticación multifactor, MFA) exige evidencias de distintas factor category, por ejemplo possession e inherence; \term{WebAuthn/passkey} usa credenciales de clave pública y puede ofrecer phishing-resistant authentication. Que la autenticación tenga éxito no significa que el usuario pueda acceder a todos los recursos: la authorization todavía debe revisar group, role, attribute, device, risk y requested action.

\lecturefigure{03-identity-lifecycle.png}{La identity infrastructure abarca cuatro tipos de problemas: autenticación, autorización, aprovisionamiento/revocación, y session/audit.}{Conceptos de Okta/OIDC/SCIM/WebAuthn; redibujo conceptual.}

\begin{knowledgebox}{Lectura de la figura: Login y lifecycle operan en escalas de tiempo distintas}
La authentication responde en segundos quién es el sujeto actual; la authorization restringe los permisos en cada acción sobre un recurso; el \term{provisioning} (aprovisionamiento) crea cuentas y entitlement en las aplicaciones posteriores; \term{SCIM} (System for Cross-domain Identity Management) sincroniza usuarios y group mediante una API estándar; session/audit, en cambio, se extiende hasta el cierre de sesión, un cambio de riesgo, la baja o una investigación. El SSO puede funcionar con normalidad y aun así una falla de offboarding deja una orphan account; por eso, los indicadores de seguridad deben cubrir todo el lifecycle.
\end{knowledgebox}

\begin{warningbox}{La MFA no es la solución universal contra el session theft}
La MFA reduce de forma notable el riesgo de que alguien se apodere de una cuenta solo con la contraseña, pero si el atacante roba una bearer session ya establecida, induce al usuario a aprobar un challenge equivocado o controla el proceso de recuperación o de help-desk, todavía puede eludir las garantías esperadas. El sistema necesita phishing-resistant authenticator, session binding, risk re-evaluation, recovery control y revoke rápido, en lugar de limitarse a contar el MFA enrollment.
\end{warningbox}

\subsection{Resumen de la sección}

Lo que unifica el identity control plane es la interfaz de estado y de decisión, no la eliminación de toda la source ownership. La autenticación, la autorización, el provisioning, la session y el audit pertenecen a ciclos de vida distintos, y para cada uno hay que definir por separado la correctness y la recovery.
\section{Front-Door Reliability: seguridad y disponibilidad deben cumplirse a la vez}

El servicio de identidad funciona como puerta de entrada de la organización: si deja de estar disponible, el correo, el código, el CRM, la consola de la nube y las herramientas internas pueden fallar al mismo tiempo; si concede acceso por error, un atacante puede propagarse entre aplicaciones. Por eso la reliability no puede medirse solo con el uptime, ni la security solo con los blocked attacks. El sistema necesita un SLO conjunto de availability, latency, policy correctness, state freshness y safe change.

\subsection{SLO, blast radius y safe change}

Un \term{service-level objective} (objetivo de nivel de servicio, SLO) es una meta medible que el sistema se compromete a alcanzar dentro de una ventana estadística, por ejemplo, la proporción de autenticaciones exitosas o la p99 latency. El \term{blast radius} (radio de impacto de una falla) indica el alcance de tenants, regions, applications o principals afectados por un solo error. Una plataforma de identidad multi-tenant aumenta la eficiencia al compartir recursos, pero también exige mayor isolation, canary, rollback y dependency budget.

\lecturefigure{04-front-door-reliability.png}{La identidad es una puerta de entrada con un gran radio de impacto de fallas; hay que gestionar a la vez availability, latency, correctness y change safety.}{Subtítulos locales 04:00--06:10, 14:40--16:40; redibujo conceptual.}

\begin{knowledgebox}{Lectura de la figura: las cuatro dimensiones no se sustituyen entre sí}
La availability indica si las solicitudes legítimas se completan; la latency, si la redirect, el factor y la policy chain terminan dentro de una latencia de cola aceptable; la correctness, si el allow/deny usa el principal, la policy y la session correctos; la change safety, si una regla nueva, un connector o un cambio de código pueden validarse a pequeña escala y revertirse con rapidez. Un uptime de 99.99\% no puede ocultar una autorización errónea, y una policy extremadamente estricta tampoco justifica que nadie pueda iniciar sesión.
\end{knowledgebox}

\begin{importantbox}{El error budget debe distinguir los errores de seguridad}
Un SLO común puede escribirse como
\[
\mathrm{availability}=1-\frac{N_{failed}}{N_{eligible}}.
\]
Aquí, $N_{failed}$ son las solicitudes que cumplen las condiciones estadísticas pero no tuvieron éxito, y $N_{eligible}$ son todas las eligible requests. Sin embargo, un sistema de identidad también debe rastrear por separado los false allow, los false deny, las stale entitlements y el revoke delay; estos correctness/security errors no deben quedar diluidos en un solo número global de availability.
\end{importantbox}

\subsection{Cada acceso es una context-aware policy decision}

La \term{assurance} (nivel de garantía) describe el grado de confianza del sistema en que «el sujeto actual es realmente el titular de la cuenta», y depende del authenticator, del dispositivo y del contexto. La idea común de Okta Identity Engine y de Zero Trust es que, según la aplicación de destino y el riesgo de la solicitud, la assurance requerida también cambia: una aplicación de bajo riesgo puede reutilizar la session, mientras que una operación muy sensible exige step-up authentication.

\lecturefigure{05-access-request-path.png}{La solicitud de acceso pasa por identity/context, policy, assurance y session outcome, en lugar de limitarse a verificar la contraseña.}{Okta Identity Engine policy model; NIST SP 800-207; redibujo conceptual.}

\begin{knowledgebox}{Lectura de la figura: las decisiones de policy deben poder reproducirse}
La request lleva el principal, el device/network y la target app; el context reúne profile, group, risk y session history; la policy produce allow, deny o step-up; el outcome crea la session/token y registra un audit event. Para investigar un allow anómalo, el sistema debe conservar la policy/version, los attributes clave, la risk signal, el authenticator, la decision reason y el ID de token/session, y no solo un registro de «inicio de sesión exitoso».
\end{knowledgebox}

\subsection{Resumen de la sección}

La identity reliability significa que «el sujeto correcto accede al recurso correcto con la garantía correcta, y ese acceso puede revocarse cuando cambia el riesgo». El SLO, el radio de impacto de fallas, la policy correctness y el safe rollout deben diseñarse en conjunto.
\section{Founder Adaptation: de una organización con certezas a un sistema de baja probabilidad}

Después, la clase pasa al papel del fundador. Cuando McKinnon dejó Salesforce, perdió el equipo, la marca, el apoyo administrativo y una trayectoria de ascenso muy probable; las entradas del sistema emprendedor pasaron a ser un runway limitado, un buyer signal difuso y unas hipótesis que cambian todo el tiempo. Esta sección conserva esa teacher voice porque explica cómo una empresa de infraestructura, en sus primeras etapas, puede lograr avances verificables aun sin contar con evidencia completa.

\subsection{Resource Reset y daily evidence}

El fundador no puede usar la «convicción» como pretexto para rechazar los datos. Un ciclo más sólido consiste en descomponer la gran visión en evidence que pueda actualizarse cada día: quién está dispuesto a dar una entrevista, qué pain desencadena un presupuesto, si el prototype se conecta a un directory real, si el first tenant es estable, si el primer pago se repite. Cada paso reduce un tipo de incertidumbre y da fundamento a la siguiente ronda de contrataciones, de financiamiento o de decisiones de producto.

\lecturefigure{06-founder-adaptation.png}{Pasar de una empresa consolidada a un equipo emprendedor significa que los recursos y la certeza vuelven a cero; después, el sistema de ejecución se reconstruye con daily evidence.}{Subtítulos locales 06:00--12:20; concepto redibujado.}

\begin{knowledgebox}{Lectura de la figura: Belief y probability pueden coexistir}
Resource reset indica que las ventajas de la organización anterior no se trasladan de forma automática; daily evidence convierte la visión en señales de customer/prototype; low odds exige definir con claridad el runway y la stopping rule; team belief, por su parte, requiere rumbo, honestidad y avances visibles. Lo que dice el ponente, «hay que creer incluso cuando no crees», no debe entenderse como engañar al equipo, sino como la capacidad del líder de organizar la acción aun reconociendo que las probabilidades están en contra.
\end{knowledgebox}

\teachervoice{McKinnon reconoce con franqueza que su motivación para irse fue tanto la oportunidad técnica como el deseo personal de «ser su propio jefe»; lo verdaderamente difícil fue tomar una decisión en común frente a la crisis financiera, un recién nacido y los riesgos para la familia. La clase saca el founder risk del relato del héroe individual y lo devuelve a la familia, los recursos y el costo de oportunidad.}

\begin{warningbox}{No conviertas el survivorship bias en una fórmula para emprender}
El éxito posterior de Okta no demuestra que «dejar una gran empresa», «perseverar lo suficiente» o «la experiencia de la madurez» produzcan resultados por necesidad. Los principios transferibles son elegir un market acorde con las propias capacidades, obtener buyer evidence lo antes posible, limitar los riesgos irreversibles y actualizar el plan cuando aparezca nueva evidencia; el resultado sigue dependiendo del momento, la competencia, el capital y la suerte.
\end{warningbox}

\subsection{Resumen de la sección}

Founder adaptation es pasar de los inherited resources a un evidence loop. La convicción sostiene la acción continua; la probabilidad y las métricas se encargan de corregir el rumbo. Si falta cualquiera de las dos, la organización pierde el sentido de la realidad.
\section{CEO y Board: construir un circuito de información sin distorsión}

Cuando la empresa crece, el CEO deja de resolver directamente todos los problemas y se convierte en el punto donde convergen la información, el capital, el talento y el riesgo. La «soledad» del cargo no suele deberse a que nadie le hable. Se debe a que las malas noticias se filtran al pasar por la jerarquía, el equipo espera respuestas definitivas y el consejo de administración solo puede juzgar con la evidence que le proporciona la dirección. La solución sistémica consiste en definir con claridad los decision rights y un unfiltered information loop.

\subsection{Bad news, options y accountability}

El equipo directivo debe convertir la operating truth en algo que se pueda discutir: qué ocurrió con los indicadores, dónde están los riesgos para los clientes o de seguridad, qué root cause ya se comprobaron y cuáles son solo una hypothesis, y qué options, trade-off y time constraint existen. La función del consejo de administración no es dirigir la empresa en lugar de la dirección. Su función es cuestionar los supuestos, asignar el capital y supervisar los riesgos importantes. El CEO, por su parte, responde por la decisión final y por el follow-through.

\lecturefigure{07-ceo-board-loop.png}{El CEO y el consejo de administración necesitan un circuito sin distorsión que vaya de la operating truth a las options, la decision y el follow-through.}{Subtítulos locales 12:00--14:40; redibujo conceptual.}

\begin{knowledgebox}{Lectura de la figura: el Dashboard no sustituye el relato de las malas noticias}
Las Metrics le indican al consejo «dónde hay una anomalía»; la customer evidence y la incident timeline explican «por qué importa»; las options precisan «qué se puede hacer», y el owner/milestone permite revisar la decisión después. Si el CEO solo informa de los problemas ya resueltos, el board conocerá el contexto por primera vez justo cuando de verdad se necesite su apoyo. Si el board pasa por encima del CEO y dirige la ejecución directamente, rompe la cadena de responsabilidad.
\end{knowledgebox}

\teachervoice{Al recordar sus primeras interacciones con el consejo de administración, el ponente considera que se esforzó demasiado por demostrar que «todo estaba bajo control». Más adelante encontró un método más eficaz: llevar las dificultades y las malas noticias a la discusión con anticipación. Para un equipo de frontier systems, esto equivale a hacer que los riesgos aparezcan cuando todavía se puede elegir, en lugar de escalarlos solo después de un outage o una breach.}

\subsection{Resumen de la sección}

La calidad del CEO/board loop depende de que la información esté completa, de la option clarity y de los límites de responsabilidad. La soledad no se resuelve con más status meeting. Se resuelve con una estructura de gobierno capaz de dar cabida a la incertidumbre y a las malas noticias.
\section{Security Incident: de la crisis a la mejora verificable}

Un identity provider es un objetivo de alto valor, y un incidente de seguridad desafía directamente los supuestos de confianza más fundamentales del cliente. La clase enfatiza el dolor que trae el incidente y la transformación cultural; el RCA oficial ofrece una perspectiva de sistema más concreta: cómo se expuso la credential, por qué la consulta de registros pasó eventos por alto, cómo se aprovechó el session token, cuándo se hizo la containment, cómo se notificó a los clientes y qué control se modificaron de forma permanente.

\subsection{Las cinco etapas de la incident response}

Un incident lifecycle completo incluye, como mínimo, detection, containment, investigation, remediation y notification. La detection puede provenir de la telemetry interna o de clientes y socios; la containment debe deshabilitar la account, hacer revoke de session/token y restringir las interfaces; la investigation establece el timeline y el affected scope; la remediation corrige credential, logging, workflow y product control; la notification, por su parte, entrega a los clientes afectados información sobre la que puedan actuar.

\lecturefigure{08-incident-response.png}{Un incidente de seguridad solo se da por terminado cuando se cierra el ciclo de aprendizaje de detection, containment, investigation, remediation y notification.}{Subtítulos locales 15:20--19:10; RCA del support-system de Okta de 2023; rediseño conceptual.}

\begin{knowledgebox}{Lectura de la figura: el RCA debe explicar «por qué no se detectó antes»}
La root cause no solo busca el error inicial; también debe explicar por qué las defensas no lo impidieron, por qué la telemetry no lo mostró, por qué la investigation query lo pasó por alto, por qué se retrasó el session revoke y cómo entraron las señales de los clientes en la respuesta. El RCA oficial de 2023 explica que distintas rutas de acceso a archivos generan distintos log event, por lo que la consulta inicial quedó incompleta; este detalle muestra que el propio esquema de observabilidad también forma parte de los security control.
\end{knowledgebox}

\begin{importantbox}{La containment y la remediation son distintas}
Deshabilitar la service account, hacer revoke de la session y bloquear direcciones IP es containment: su objetivo es detener el daño en curso. Impedir que los profile personales guarden credential de la empresa, reforzar la telemetry de acceso a archivos, vincular las session de administrador y mejorar la customer notification es remediation: su objetivo es reducir la probabilidad de recurrencia y acortar el detection time futuro. Ambas requieren distintos owner y distintas evidencias de cumplimiento.
\end{importantbox}

\begin{warningbox}{La transparencia no es publicar una entrada de blog}
La customer trust requiere un scope oportuno, los objetos afectados, la recommended action, lo que se sabe que se desconoce y actualizaciones posteriores. Dar demasiado pronto conclusiones definitivas no comprobadas se vuelve contra la confianza, y divulgar demasiado tarde le quita al cliente tiempo de containment. La incident communication debe ir a la par de la forensic confidence y conservar un decision log de quién sabía qué y en qué momento.
\end{warningbox}

\teachervoice{McKinnon dice que un incidente de seguridad es una de las etapas más difíciles de una carrera profesional, porque el headline asocia directamente a la empresa de identidad con «fue vulnerada». El juicio de la clase que vale la pena conservar es este: la confianza no se recupera con consignas; solo puede recuperarse convirtiendo el dolor en recursos, conductas y controles verificables orientados a security-first.}

\subsection{La security-first culture es un operating system}

La culture no es una value list colgada en la pared, sino lo que los líderes preguntan una y otra vez, lo que están dispuestos a retrasar en un launch, a quién asignan presupuesto, qué recompensan los ascensos y si las malas noticias se castigan. Si «security first» solo exige que el security team trabaje horas extra, mientras el product roadmap, los sales commitment y la executive review no cambian, la señal real que recibe la organización sigue siendo feature first.

\lecturefigure{09-security-first-culture.png}{La security-first culture se logra con la conducta de los líderes, los incentivos, los recursos y la evidencia en conjunto.}{Subtítulos locales 18:00--21:30; rediseño conceptual.}

\begin{knowledgebox}{Lectura de la figura: la culture puede auditarse}
En la leader behavior se observa si la alta dirección pregunta por los riesgos por iniciativa propia y acepta la escalation; en el incentive, los objetivos, el desempeño y los launch criteria; en el resource, si la renovación de la plataforma, el security staffing y el customer support reciben presupuesto; en la evidence, si los control test, los incident metric, las audit y los postmortem entran en la review. Aunque la culture es difícil de cuantificar con precisión, puede observarse a través de conductas repetidas y de la resource allocation.
\end{knowledgebox}

\subsection{Resumen de la sección}

El objetivo de la incident response no es «volver a la luz verde», sino establecer un learning loop que pueda revisarse. Security-first solo surte efecto en la siguiente situación de presión si se incorpora a las decisiones, a los recursos y a la evidencia de ingeniería.
\section{Attack Asymmetry, Shared Signals y Zero Trust}

Al atacante le basta con encontrar una credential, una session, un help-desk o un integration gap; el equipo defensor, en cambio, debe coordinar endpoint, identity, application, data y human workflow. En clase se vinculó esta asimetría con los obstáculos para compartir información: el diseño de sistemas necesita defense in depth y también intercambiar risk signal dentro de los límites de la privacidad y la ley, para reducir el tiempo que cada organización tarda en descubrir por su cuenta el mismo patrón de ataque.

\subsection{Una sola debilidad frente a varias capas de defensa}

El equipo defensor no puede cargar todo el riesgo sobre el IdP. Endpoint compromise, phishing, malware, OAuth consent, stale account, over-privileged service account y application bug pueden, cada uno, eludir una verificación única. La ventaja de la identity layer es que ve el inicio de sesión entre aplicaciones y el estado de factor, session y group, pero aun así necesita recibir indicadores de EDR, network, application y customer-reported indicators.

\lecturefigure{10-shared-defense.png}{El atacante necesita un solo hueco; el defensor debe coordinar controles en varias capas y señales entre organizaciones.}{Subtítulos locales 21:20--24:10; redibujo conceptual.}

\begin{knowledgebox}{Lectura de la figura: la shared signal también debe minimizarse}
Un indicator puede ser una suspicious IP, session risk, credential compromise o device posture; antes de compartirlo deben definirse schema, confidence, purpose, retention y recipient. Si se transmite muy poco, se pierde valor defensivo; si se transmite demasiado, aumenta el riesgo para datos privados y comercialmente sensibles. La signal debe activar una policy evaluation local, no otorgar a la fuente externa la facultad de emitir directamente el deny final.
\end{knowledgebox}

\subsection{Zero Trust: confianza explícita, contextual y revocable}

\term{Zero Trust} no significa «nunca confiar en nadie», sino no conceder confianza permanente por la ubicación en la red o por un solo inicio de sesión. NIST SP 800-207 enfatiza que cada acceso se decida de forma explícita según subject, asset, environment y policy, que se aplique least privilege y que se suponga que la red pudo haber sido comprometida. Las identity policy modernas, además, reciben señales de riesgo de forma continua durante la session y aplican step-up, reducción de privilegios o logout.

\lecturefigure{11-zero-trust-policy.png}{Zero Trust convierte el contexto de identidad, dispositivo, riesgo y recurso en una policy continua y revocable.}{NIST SP 800-207; Okta policy/Identity Threat Protection; redibujo conceptual.}

\begin{knowledgebox}{Lectura de la figura: la policy evaluation no ocurre solo al iniciar sesión}
Identify establece el principal y el authenticator; evaluate reúne device, network y risk; enforce produce least privilege, step-up o deny; re-evaluate aplica revoke o universal logout cuando cambia el riesgo de la session. La evaluación continua no implica vigilancia ilimitada, sino elegir las señales necesarias para cada action de alto valor y dar a la session una expiry y una revocation path claras.
\end{knowledgebox}

\begin{warningbox}{Zero Trust no corrige un entitlement desordenado}
Si la organización no sabe quién es dueño de cada service account, qué permisos corresponden a cada group, cuándo se propaga una baja de personal o si una aplicación admite revoke, autenticar con más frecuencia solo sirve para verificar una y otra vez al mismo sujeto con privilegios excesivos. Zero Trust presupone inventory, ownership, least privilege y lifecycle hygiene.
\end{warningbox}

\subsection{Resumen de la sección}

La attack asymmetry exige defensa en varias capas y señales compartidas; Zero Trust, a su vez, convierte esas señales en decisiones de acceso explícitas, de mínimo privilegio y revocables. Ambas dependen de un identity state limpio y de límites de auditoría claros.
\section{AI Agent Identity: Delegation en lugar de secretos compartidos}

La clase ya señalaba en 2025 la brecha de seguridad de los agent frameworks: los scripts y las service accounts llevan tiempo ejecutando tareas en nombre de personas o sistemas; el LLM agent solo hace que el comportamiento sea más dinámico, que haya más herramientas y que la intención expresada en lenguaje natural sea más difícil de acotar con precisión. El problema central no es «si el agent tiene nombre de usuario», sino quién lo autoriza, en nombre de quién actúa, a qué audience puede acceder, qué actions puede ejecutar, cuándo caduca y cómo se revoca.

\subsection{Principal, delegation y OAuth}

\term{service account} es una cuenta no humana que usan el software o una workload; \term{non-human identity} (NHI) designa en general a un principal de tipo service, workload, device o agent. \term{OAuth 2.0} es un delegated authorization framework que permite al client obtener un scoped access token sin tener que poseer la contraseña del usuario; \term{OpenID Connect} (OIDC) proporciona una capa de identidad sobre OAuth y expresa el resultado de la autenticación mediante un ID token.

\lecturefigure{12-agent-delegation.png}{El agent necesita un principal propio y una delegation acotada, en lugar de heredar todos los secretos del usuario o de la máquina.}{Subtítulos locales 24:10--27:50; OAuth token exchange; redibujo conceptual.}

\begin{knowledgebox}{Lectura de la figura: distinguir «quién es el agent» de «en nombre de quién actúa»}
El user/owner aporta la intent y el consent; el agent principal acredita la workload identity; el authorization server emite el token con base en la delegation, el scope, la audience y la expiry; después, el tool/resource server realiza el policy check y escribe el audit. \term{token exchange} permite al sistema intercambiar un token existente por uno nuevo dirigido a una audience específica y con permisos más reducidos, conservando la relación entre actor y subject, para no entregar directamente al agent el broad token del usuario.
\end{knowledgebox}

\teachervoice{McKinnon advierte a los estudiantes que no se dejen confundir por la idea de que «el agent es algo completamente nuevo»: los cron jobs, los scripts y las service accounts llevan mucho tiempo actuando en nombre de los sistemas. Lo verdaderamente nuevo es la ampliación de las capacidades y de la autonomía, mientras que muchos frameworks todavía dejan la security para el final e incluso permiten que las máquinas de desarrollo acumulen tokens con acceso a todos los sistemas.}

\subsection{El token es una capability con ciclo de vida}

\term{access token} es la credencial que expresa una autorización ante el resource server; un bearer token puede usarlo cualquiera que lo obtenga, por lo que es necesario limitar su scope, su audience y su lifetime, y adoptar en lo posible tokens sender-constrained/bound. Cuando la agent task termina, el owner deja la organización, el device presenta riesgo, el modelo se deshabilita o cambia la tool policy, debe ser posible hacer revoke de la session/token existente.

\lecturefigure{13-token-lifecycle.png}{El token es una capacidad de acceso temporal que requiere un ciclo de vida completo: issue, bind, use y revoke/retire.}{Conceptos de OAuth/OIDC/WebAuthn; redibujo conceptual.}

\begin{knowledgebox}{Lectura de la figura: el least privilege debe concretarse por tarea}
En la fase de issue se confirman el principal, la audience, el scope y la expiry; en la fase de bind se asocia el token a una key/device/session, lo que reduce su utilidad tras un token theft; en la fase de use, el resource server verifica la action y el trace ID; en la fase de revoke/retire se responde a risk, job completion u owner removal. Si un agent solo necesita leer un issue y crear un borrador, no debe recibir repo admin, billing ni un refresh token de larga duración.
\end{knowledgebox}

\begin{importantbox}{Campos mínimos de un delegation record}
Cada delegation a un agent debe registrar como mínimo: human/organizational owner, agent principal, subject, requested task, resource audience, approved scopes, policy/version, issue/expiry, token/session ID, tool actions, result y revocation reason. Solo así el incident response puede responder «qué persona autorizó a qué agent a hacer qué en qué sistema».
\end{importantbox}

\begin{warningbox}{La agent memory no debe convertirse en un credential warehouse}
El prompt, el memory store, la trace y el local workspace suelen diseñarse como contexto recuperable; si en ellos se mezclan API keys, refresh tokens o session cookies, las salidas del modelo, los complementos y los registros de depuración pueden ampliar la superficie de filtración. Los secrets deben guardarse en un vault dedicado, inyectarse por task mediante un short-lived credential broker y excluirse del contexto del modelo y de los registros ordinarios.
\end{warningbox}

\subsection{Resumen de la sección}

Lo esencial de la agent identity es la separación entre actor y subject, la delegation, el least privilege, el short-lived token y el audit. Compartir la contraseña del usuario o un token de larga duración convierte un solo compromise de un agent o de una máquina en un acceso persistente a varios sistemas.
\section{Auth0 Acquisition: Primitive compartidas, sin prisa por fusionar todos los límites}

McKinnon describe a Okta y Auth0 como complementarias: la primera se ha orientado desde hace tiempo a workforce/IT, y la segunda es conocida por su customer identity developer-first. Ambas comparten las primitive de directory, authentication, authorization, protocol, security y scale, pero atienden a distintos buyer, end user, SDK, customization y release expectation. El primer principio de la arquitectura de una adquisición no es dibujar cuanto antes un solo org chart, sino proteger el customer contract.

\subsection{Workforce Identity y Customer Identity}

La workforce identity se ocupa de empleados, contractor, IT administration, enterprise app y lifecycle; la customer identity se ocupa de los usuarios finales del producto, developer integration, custom UX, social login, application-specific profile y tráfico de internet. Una plataforma común puede reutilizar protocol, threat intelligence y reliability practice, pero la runtime migration, el tenant model y la developer experience requieren validación independiente.

\lecturefigure{14-two-platform-acquisition.png}{Una adquisición debe respetar los límites de usuarios, flujos de trabajo e infraestructura que distinguen a la workforce identity de la customer identity.}{Subtítulos locales 30:00--36:20; materiales oficiales de la adquisición de Okta/Auth0 de 2021; redibujo conceptual.}

\begin{knowledgebox}{Lectura de la figura: una shared primitive no equivale a un shared runtime}
Ambos tipos de producto necesitan directory, OAuth/OIDC, MFA, security y scale; pero el workforce buyer puede exigir IT policy, employee lifecycle y enterprise integration, mientras que la customer identity pone el acento en SDK, experiencias con la marca propia, picos extremos y developer control. Después de la adquisición conviene unificar primero la threat intelligence, la support escalation y la strategic interface, y luego decidir con evidencia de compatibility, SLO y migration si se fusiona el runtime.
\end{knowledgebox}

\teachervoice{Al recordar la adquisición, el ponente considera que el escepticismo de los inversionistas hacia las M\&A fue mayor de lo que él esperaba; también reconoce que, en un entorno de tasas de interés bajas y crecimiento acelerado, parte de la integración del negocio pudo avanzar demasiado rápido. Esta reflexión advierte a los equipos de sistemas que el crecimiento de corto plazo oculta el integration cost, y que un viento macroeconómico a favor no debe confundirse con una arquitectura ya validada.}

\begin{warningbox}{Tres tipos de costos ocultos de la platform consolidation}
Unificar el empaquetado comercial puede ocurrir antes de que entitlement y billing sean coherentes; unificar la console puede ocultar distintos tenant/security model; unificar el backend amplía el blast radius, obliga a migrar estado y modifica los límites de cumplimiento normativo del cliente. Cada paso requiere rollback, data migration validation, customer communication y parallel run, en lugar de sustituir la aceptación técnica por «one company».
\end{warningbox}

\subsection{Resumen de la sección}

La integración posterior a una adquisición es una migración conjunta del product contract, del runtime state y de los incentivos de la organización. Compartir identity primitive puede generar sinergias, pero preservar los límites (preserve boundary) suele ser más seguro que unificar de inmediato.
\section{Durable Trust: el efecto compuesto de Reliability, Security y la recuperación transparente}

Los primeros clientes de una identity startup tienen que confiar en que un equipo pequeño puede proteger la puerta de entrada de la organización, sin paros frecuentes por mantenimiento, y en que asumirá su responsabilidad cuando ocurra un incidente. La experiencia de McKinnon en Salesforce aporta la credibility inicial, pero la confianza a largo plazo solo se acumula con production evidence: servicio estable, change correctos, incident response honesta, customer guidance e inversión continua en la plataforma.

\subsection{Trust Flywheel}

Un servicio confiable lleva a los clientes a ampliar su deployment. Más aplicaciones y más usuarios aportan una integration/context más rica, pero también amplían la responsabilidad. La security evidence y un RCA transparente ayudan al cliente a evaluar el residual risk. A su vez, la confianza del cliente trae renovaciones, reference y una adopción más profunda de la plataforma. El volante también puede girar en sentido contrario: un poorly handled incident no solo afecta el SLA del periodo, sino que lleva al cliente a reducir el alcance de su integración, lo que debilita el context futuro y el valor del producto.

\lecturefigure{15-trust-flywheel.png}{La adopción de una plataforma de identidad genera un efecto compuesto mediante la confiabilidad, la evidencia de seguridad, la confianza del cliente y la profundidad de la plataforma.}{Subtítulos locales 27:30--35:20; redibujo conceptual.}

\begin{knowledgebox}{Lectura de la figura: Trust es evidencia acumulada, no una promesa de cero incidentes}
El reliable service demuestra que se puede depender de la front door; la security evidence muestra el control y la capacidad de aprendizaje; la customer trust se refleja en implementaciones más amplias, reference y renewal; y la platform depth, a su vez, mejora la integration y el policy context. Ningún sistema grande puede prometer con honestidad que nunca tendrá incidentes, pero sí puede comprometerse a detectarlos rápido, limitar el blast radius, publicar hechos sobre los que se pueda actuar y completar una remediation verificable.
\end{knowledgebox}

\teachervoice{El ponente explicó que los primeros clientes no solo compran una funcionalidad, también evalúan «si esta persona se irá en cuanto cierre la venta y si es capaz de mantener estable el servicio». Su experiencia previa en ingeniería le dio la primera ronda de confianza, pero cada outage, breach y recuperación posterior vuelve a fijar el valor de esa confianza.}

\subsection{Resumen de la sección}

La identity trust se construye con evidencia operativa de largo plazo. Reliability, security, transparency y customer outcome se refuerzan entre sí; el eslogan de marca solo puede amplificar la evidencia que ya existe, no sustituirla.
\section{AI Adoption: del Mr. T Demo al Category-Defining Workflow}

Al final de la entrevista, McKinnon compara la IA con las primeras aplicaciones del iPhone: la capacidad del hardware o del modelo importa, sin duda, pero al principio muchas aplicaciones solo envolvían la nueva capacidad como novelty; lo que de verdad transformó la industria fue el workflow al estilo de Uber, que llegó después. Este juicio no afirma que la IA carezca de valor, sino que advierte a las empresas incumbent y a los equipos emprendedores que distingan entre capability shock, demo, repeatable workflow y organization redesign.

\subsection{Capacidad técnica e inercia organizacional}

Las grandes empresas tienen clientes, datos y procesos, pero también role, budget, support model e inercia cultural. Aunque un AI agent pueda completar parte de una tarea de support, operations o security, la organización todavía necesita reescribir accountability, quality review, exception handling, identity delegation y performance metric. Si solo se coloca un cuadro de chat delante del proceso anterior, el sistema puede terminar con una «Mr. T app» sin haber cambiado la outcome economics.

\lecturefigure{16-ai-adoption-curve.png}{Por lo general, la IA pasa primero por un capability shock y por novelty apps; solo después aparecen el category-defining workflow y el rediseño organizacional.}{Subtítulos locales 36:00--39:02; redibujo conceptual.}

\begin{knowledgebox}{Lectura de la figura: cuatro preguntas para identificar un killer workflow}
¿Completa una tarea que antes no podía realizarse de forma económica, en lugar de solo cambiar la UI? ¿Su uso genera compounding data/feedback? ¿Las fallas tienen un owner y una recovery claros? ¿La organización está dispuesta a cambiar role, process, control y pricing? Si todas las respuestas son negativas, la aplicación probablemente siga siendo una demo interesante; solo si cambian a la vez el workflow, la distribution y la governance puede la tecnología formar una nueva category.
\end{knowledgebox}

\teachervoice{McKinnon no adopta un relato de AI doom ni afirma que el killer use case ya esté definido. Compara la etapa actual con la aplicación de Mr. T de los primeros tiempos del iPhone, que decía frases ingeniosas: todos saben que la plataforma importa, pero «un workflow completamente nuevo que nace de una nueva capacidad, como Uber» todavía está en formación.}

\begin{warningbox}{La incumbent inertia es a la vez un riesgo y un margen de seguridad}
Los procesos anteriores frenan la automatización valiosa, pero también pueden impedir por un tiempo que un agent sin gobernanza obtenga permisos de producción. El rediseño organizacional debe, al mismo tiempo, eliminar los handoff ineficaces y añadir delegation, review, audit y fallback; «AI-first» no puede servir de pretexto para eludir identity, change management e incident response.
\end{warningbox}

\subsection{Resumen de la sección}

El valor de la IA va de la capability al workflow y, después, a la organization. En este trayecto, la identity infrastructure no es un componente de inicio de sesión que se añade al final, sino la condición previa para la agent accountability, el least privilege y el cross-tool audit.
\section{Síntesis y ampliación}

Esta clase extrae de la experiencia de fundación y operación de Okta principios transferibles de la identity infrastructure:

\begin{enumerate}
  \item \textbf{La transition necesita un wedge}: el cambio tecnológico abre una ventana; el acceso estrecho por el que el cliente está dispuesto a pagar hoy se encarga del arranque, y la profundidad de la plataforma, del valor a largo plazo;
  \item \textbf{Identity es el control plane}: directory, authentication, authorization, provisioning, session y audit determinan en conjunto el acceso;
  \item \textbf{La front door exige a la vez seguridad y disponibilidad}: availability, latency, correctness, blast radius y safe change deben validarse en conjunto;
  \item \textbf{La convicción debe conectarse con la evidence}: un entorno de baja probabilidad requiere daily progress, buyer signal y una uncertainty honesta, no rechazar los datos;
  \item \textbf{La gobernanza depende de información sin distorsión}: el CEO y el board deben convertir las malas noticias en options, decision rights y follow-through;
  \item \textbf{Un incident debe cerrar el ciclo de aprendizaje}: containment, RCA, remediation, notification y culture change son todos indispensables;
  \item \textbf{Zero Trust establece una confianza revocable}: el contexto, least privilege, step-up y continuous re-evaluation dependen de un lifecycle limpio;
  \item \textbf{El agent usa delegation}: un principal independiente, actor/subject, short-lived scoped token, vault y audit sustituyen a los secretos compartidos;
  \item \textbf{En una adquisición, primero se protege el customer contract}: compartir primitives no significa fusionar de inmediato el runtime, los equipos y los flujos de trabajo de producto;
  \item \textbf{La IA pasa de la demo al workflow}: el category value requiere que cambien en conjunto las capacidades, los procesos, la organización y los controles de seguridad.
\end{enumerate}

\begin{importantbox}{Tarea de Identity Architecture}
Entrega un diseño de una página para un sistema en el que «un AI agent lee tickets en nombre de un empleado y crea borradores de código»: dibuja al human, el agent principal, el authorization server, el ticket system, el code host y el audit; define el flujo OIDC/OAuth/token-exchange, scope, audience, expiry, vault, session revoke y owner offboarding; después escribe las estrategias de degradación ante IdP outage, stale group, token theft y agent compromise. Queda prohibido que el agent herede directamente los token de larga duración del usuario.
\end{importantbox}

\subsection{Lecturas adicionales}

{\small
\begin{itemize}[nosep,leftmargin=*]
  \item \href{https://developer.okta.com/docs/concepts/universal-directory/}{\textbf{Okta Universal Directory}}: identity data model unificado y source integration.
  \item \href{https://help.okta.com/oie/en-us/content/topics/identity-engine/policies/about-policies.htm}{\textbf{Okta Identity Engine policies}}: global session y app sign-in assurance.
  \item \href{https://sec.okta.com/articles/2023/11/unauthorized-access-oktas-support-case-management-system-root-cause/}{\textbf{Okta support-system incident RCA}}: timeline, logging gap y remediation.
  \item \href{https://csrc.nist.gov/pubs/sp/800/207/final}{\textbf{NIST SP 800-207}}: Zero Trust Architecture.
  \item \href{https://www.rfc-editor.org/rfc/rfc8693}{\textbf{RFC 8693}}: OAuth 2.0 Token Exchange.
  \item \href{https://www.rfc-editor.org/rfc/rfc7644}{\textbf{RFC 7644}}: SCIM provisioning protocol.
  \item \href{https://www.w3.org/TR/webauthn-3/}{\textbf{WebAuthn Level 3}}: phishing-resistant public-key authentication.
\end{itemize}
}

\end{document}
